\section{Сценарии тестирования}\label{sec:scenarios}

Набор сценариев был выбран так, чтобы разделить четыре различные причины
деградации качества: обычную стохастичность потока задач, кратковременную
перегрузку, гетерогенность вычислительных ресурсов и структурные отказы
узлов. Во всех сценариях использовалась одна и та же базовая конфигурация
сети, описанная в главе~5; менялись только параметры поступления задач и
режимы доступности ресурсов.

\subsection{Сценарий 1: стабильная нагрузка}

Базовый сценарий служит контрольной точкой для всего сравнения. Интенсивность
поступления задач задаётся параметром $\lambda = 2.5$, всплески нагрузки и
отказы узлов отключены, длительность эпизода составляет $80$ шагов.

Цель сценария --- проверить, как методы ведут себя в спокойной стационарной
среде, где нет резких внешних шоков и на первый план выходят свойства самой
политики распределения: эффективность, admission control, задержка и
справедливость распределения нагрузки.

\subsection{Сценарий 2: пиковая нагрузка}

Во втором сценарии моделируется нестационарный поток задач с базовой
интенсивностью $\lambda = 4.0$ и двумя детерминированными окнами перегрузки.
Внутри этих окон интенсивность входного потока умножается на коэффициент
$2.2$, а длительность эпизода увеличивается до $100$ шагов. Такой сценарий
позволяет анализировать не только средний итоговый результат, но и
поведение системы до, во время и после шокового увеличения нагрузки.

\begin{table}[htbp]
	\centering
	\caption{Интервалы всплесков нагрузки в сценарии 2}
	\label{tab:load-spikes-windows}
	\resizebox{\textwidth}{!}{\begin{tabular}{llllll}
\toprule
Метод & Интервал & SW & Принятие задач, \% & Задержка, мс & До дедлайна \\
\midrule
Heuristic-LoadAware & До пика & 5.926 & 68.222 & 176.708 & 0.873 \\
Heuristic-LoadAware & Между пиками & 5.233 & 66.185 & 200.502 & 0.947 \\
Heuristic-LoadAware & Пик & 8.828 & 65.242 & 198.807 & 0.982 \\
Heuristic-LoadAware & После пика & 3.998 & 66.376 & 176.057 & 0.861 \\
MA-VCG & До пика & 5.926 & 68.222 & 176.708 & 0.873 \\
MA-VCG & Между пиками & 5.233 & 66.185 & 200.502 & 0.947 \\
MA-VCG & Пик & 8.828 & 65.242 & 198.807 & 0.982 \\
MA-VCG & После пика & 3.998 & 66.376 & 176.057 & 0.861 \\
MA-VCG-QMIX & До пика & 5.751 & 64.802 & 177.211 & 0.875 \\
MA-VCG-QMIX & Между пиками & 5.256 & 65.163 & 201.318 & 0.940 \\
MA-VCG-QMIX & Пик & 8.526 & 63.619 & 196.180 & 0.976 \\
MA-VCG-QMIX & После пика & 3.869 & 65.485 & 182.166 & 0.871 \\
QMIX & До пика & 5.452 & 62.276 & 170.240 & 0.857 \\
QMIX & Между пиками & 4.960 & 63.434 & 196.135 & 0.921 \\
QMIX & Пик & 7.920 & 59.893 & 198.434 & 0.982 \\
QMIX & После пика & 3.717 & 61.429 & 178.957 & 0.862 \\
Random & До пика & 4.522 & 51.004 & 164.700 & 0.815 \\
Random & Между пиками & 3.909 & 51.299 & 184.443 & 0.882 \\
Random & Пик & 6.610 & 51.093 & 191.741 & 0.969 \\
Random & После пика & 2.894 & 48.886 & 163.698 & 0.789 \\
\bottomrule
\end{tabular}
}
\end{table}

\subsection{Сценарий 3: разнородные ресурсы}

В этом режиме интенсивность поступления задач составляет $\lambda = 3.0$,
но сами периферийные узлы становятся гетерогенными по CPU- и
memory-ёмкости. Тем самым возникает необходимость не только принимать или
отклонять заявки, но и учитывать специализацию узлов, избегая как
систематической недозагрузки <<сильных>> узлов, так и перегрузки тех
ресурсов, которые кажутся наиболее выгодными с локальной точки зрения.
Длительность эпизода сохраняется равной $80$ шагам.

\subsection{Сценарий 4: отказы узлов}

Четвёртый сценарий имитирует структурное нарушение среды. Интенсивность
потока равна $\lambda = 3.0$, но на $40$-м шаге эпизода из строя выводится
$20\%$ периферийных узлов, после чего их восстановление происходит через
$25$ шагов. Общая длительность эпизода составляет $90$ шагов. Такой режим
позволяет оценить глубину просадки метрик сразу после отказа, устойчивость
перераспределения нагрузки и качество восстановления до конца эпизода.

В совокупности четыре сценария покрывают как <<нормальный>> режим
эксплуатации, так и три принципиально различных стресс-режима. Поэтому
выводы по результатам нельзя свести к одному усреднённому числу: разные
методы оказываются сильнее при разных типах возмущений, и именно это
сравнение является содержательным результатом главы.
